\section{Haar measure on a compact Lie group}\label{haar:sec:haar}
On a finite group $G$, with any function $f\colon G\to\mathbb C$, the average
\begin{equation}\label{haar:finite-average}
 \int_G f\,dh:=\frac1{|G|}\sum_{V\in G}f(V)
\end{equation}
is a number. Here $|G|$ is the \emph{order} of $G$, the number of its
elements. It is a positive integer, since $G$ contains at least the
identity; Lee writes $k$ for it in the proof quoted in
Proposition~\ref{haar:prop:lee-16-10}. The average does not change when
$f(V)$ is replaced by $f(gV)$ or $f(Vg)$ for a fixed $g\in G$: the maps
$V\mapsto gV$ and $V\mapsto Vg$ are bijections of $G$, so they only permute
the $|G|$ terms of the sum. The scalar product of
Definition~\ref{dgdt:Def:2.1scalarproductG} is this average applied to
$\overline{f_1}f_2$. Haar measure is the same average on a compact Lie group
such as $U(1)$, $SU(2)$ or $SO(3)$. A lattice gauge integral is a product of
these averages, one for each link. Left and right invariance make a gauge
transformation an ordinary change of variables there.

\paragraph{Reading.} Lee constructs the invariant volume form in
Proposition~16.10 and the Remark after it (pp.~410--411)
\cite{dgdt:JLee2012}; both are quoted below. Kosmann-Schwarzbach states
existence, uniqueness and right invariance on a compact group in Chapter~3,
\S2, Theorem~2.1 and Definition~2.2 (pp.~34--35), and computes $O(2)$,
$SU(2)$ and $SO(3)$ in Exercises~3.7 and 5.4--5.6 (pp.~45, 79--80)
\cite{dgdt:YKosmann-Schwarzbach2010}. Hall uses the invariant integral to
make representations unitary in \S4.4 (pp.~93--94). He passes from the form
to a measure in Appendix~D (pp.~435--436). He states the Weyl integral formula
as Theorem~11.30 (p.~330), with the $SU(2)$ case in Example~11.33 (p.~332),
and proves orthonormality of characters as Theorem~12.15 (p.~352)
\cite{BCHall2015}. Each statement and its derivation appears in this section.

\subsection{What is being asked for}
Let $G$ be a compact Lie group of dimension $n$ and write
$L_g(V)=gV$ and $R_g(V)=Vg$.
\begin{definition}\label{haar:def:haar}
A \emph{normalized Haar measure} on $G$ is a Borel probability measure $h$
such that, for every continuous $f\colon G\to\mathbb C$ and every
$g_L,g_R\in G$,
\begin{equation}\label{haar:bi-invariance}
 \int_G f(g_LVg_R)\,h(dV)=\int_G f(V)\,h(dV),\qquad h(G)=1.
\end{equation}
\end{definition}
A finite group is a Lie group of dimension zero, and
Eq.~\eqref{haar:finite-average} is its Haar measure. For a Lie group of
positive dimension, $\int_G f\,dh$ will be the integral of $f$ times a
volume form, computed in charts as in Lee, Chapter~16. All integrands below
are continuous. The passage from an integral of continuous functions to a
Borel measure is made in Section~\ref{haar:sec:riesz}; it is needed only to
speak of the probability $h(A)$ of a set $A$.

Kosmann-Schwarzbach states the general result for locally compact groups
(Ch.~3, \S2, pp.~34--35):
\begin{theorem}[Kosmann-Schwarzbach, Theorem~2.1, p.~34 \cite{dgdt:YKosmann-Schwarzbach2010}]\label{haar:thm:KS-2-1}
Let $G$ be a locally compact group.
\begin{enumerate}
\item[(i)] There exists a positive measure on $G$, finite on compact sets
and not identically zero, that is left-invariant, i.e., for each integrable
function $f$ and for every $h\in G$,
\[
 \int_G f(hg)\,d\mu(g)=\int_G f(g)\,d\mu(g).
\]
Such a measure is unique up to multiplication by a positive real number. If
$f$ is a continuous, nonnegative function, then $\int_G f(g)\,d\mu(g)=0$,
implies $f=0$.
\item[(ii)] If $G$ is compact, there is a unique left-invariant measure $\mu$
such that $\int_G d\mu(g)=1$.
\item[(iii)] On a compact group, every left-invariant measure is
right-invariant.
\end{enumerate}
\end{theorem}
\begin{definition}[Kosmann-Schwarzbach, Definition~2.2, p.~35 \cite{dgdt:YKosmann-Schwarzbach2010}]
On a compact group, the unique left- and right-invariant measure of total
mass $1$ is called the Haar measure.
\end{definition}
In her statement $h$ is a group element; in ours $h$ is the measure.
Kosmann-Schwarzbach refers to Loomis for the proof of (i) on a general
locally compact group. For a compact \emph{Lie} group the following
theorem is proved completely in the next three subsections: a
left-invariant form (existence), its right invariance (compactness), and
uniqueness together with inversion (Fubini).
\begin{theorem}\label{haar:thm:exist-unique}
A compact Lie group has exactly one normalized Haar measure. It is also
invariant under inversion: $\int_Gf(V^{-1})\,h(dV)=\int_Gf(V)\,h(dV)$.
\end{theorem}

\subsection{Existence: a left-invariant volume form}
Choose a basis $\varepsilon^1,\dots,\varepsilon^n$ of the dual
$\mathfrak g^*=T_e^*G$ and put $\omega_e=\varepsilon^1\wedge\cdots\wedge
\varepsilon^n$. Transport it from the identity by left translation:
\begin{equation}\label{haar:omega-def}
 \omega_V:=(dL_{V^{-1}})_V^*\,\omega_e .
\end{equation}
For a matrix group, $T_VG=\{VX:X\in\mathfrak g\}$ and
$(dL_{V^{-1}})_V(VX)=V^{-1}VX=X$, so
\begin{equation}\label{haar:omega-matrix}
 \omega_V(VX_1,\dots,VX_n)=\omega_e(X_1,\dots,X_n).
\end{equation}
Equivalently, write the Maurer--Cartan form as
$\Theta=V^{-1}dV=\Theta^kX_k$ in a basis $X_k$ of $\mathfrak g$ dual to
$\varepsilon^k$. Then
\begin{equation}\label{haar:maurer-cartan}
 \omega=\Theta^1\wedge\cdots\wedge\Theta^n,\qquad
 (gV)^{-1}d(gV)=V^{-1}g^{-1}g\,dV=V^{-1}dV .
\end{equation}
The second equation shows that $\Theta$ is left invariant. This form of
$\omega$ is the one used for the Euler-angle computations of
Section~\ref{haar:sec:euler}.

To check left invariance of $\omega$ directly, use the chain rule and
$L_{(gV)^{-1}}\circ L_g=L_{V^{-1}g^{-1}g}=L_{V^{-1}}$:
\begin{align}
 (L_g^*\omega)_V&=(dL_g)_V^*\,\omega_{gV}
 =(dL_g)_V^*\,(dL_{(gV)^{-1}})_{gV}^*\,\omega_e\notag\\
 &=\big(d(L_{(gV)^{-1}}\circ L_g)\big)_V^*\,\omega_e
 =(dL_{V^{-1}})_V^*\,\omega_e=\omega_V .\label{haar:left-invariant-form}
\end{align}
Lee proves the same thing with a left-invariant frame $E_1,\dots,E_n$
(a basis of the Lie algebra) and its dual coframe $\varepsilon^i$. Left
invariance of $E_j$ gives
$(L_g^*\varepsilon^i)(E_j)=\varepsilon^i(L_{g*}E_j)=\varepsilon^i(E_j)
=\delta^i_j$, so $L_g^*\varepsilon^i=\varepsilon^i$, and then
$L_g^*(\varepsilon^1\wedge\cdots\wedge\varepsilon^n)
=L_g^*\varepsilon^1\wedge\cdots\wedge L_g^*\varepsilon^n
=\varepsilon^1\wedge\cdots\wedge\varepsilon^n$ (p.~410).

Orient $G$ by $\omega$. Lee's Proposition~15.19 says that a Lie group has
exactly two left-invariant orientations, one for each orientation of
$\mathfrak g$. The integral of a compactly supported $n$-form is invariant
under diffeomorphisms:
\begin{proposition}[Lee, Proposition~16.6(d), p.~407 \cite{dgdt:JLee2012}]\label{haar:prop:lee-16-6}
Suppose $M$ and $N$ are non-empty oriented smooth $n$-manifolds with or
without boundary, and $\omega$ is a compactly supported $n$-form on $M$.
If $F\colon N\to M$ is an orientation-preserving or orientation-reversing
diffeomorphism, then
\[
 \int_M\omega=\begin{cases}\int_NF^*\omega&\text{if $F$ is orientation-preserving,}\\
 -\int_NF^*\omega&\text{if $F$ is orientation-reversing.}\end{cases}
\]
\end{proposition}
By Eq.~\eqref{haar:left-invariant-form}, $L_g$ carries $\omega$ to itself, so
it is orientation preserving. Take $F=L_g$ and replace $\omega$ by $f\omega$:
\begin{equation}\label{haar:left-invariant-integral}
 \int_G(f\circ L_g)\,\omega
 =\int_G(f\circ L_g)\,L_g^*\omega
 =\int_GL_g^*(f\omega)
 =\int_Gf\,\omega .
\end{equation}
The first equality is Eq.~\eqref{haar:left-invariant-form}. Every
left-invariant $n$-form is $c\,\omega$ for a constant $c$, since by
Eq.~\eqref{haar:omega-def} it is determined by its value at $e$, an element
of the one-dimensional space $\Lambda^n\mathfrak g^*$. Lee's statement is:
\begin{proposition}[Lee, Proposition~16.10, p.~410 \cite{dgdt:JLee2012}]\label{haar:prop:lee-16-10}
Let $G$ be a compact Lie group endowed with a left-invariant orientation.
Then $G$ has a unique positively oriented left-invariant $n$-form
$\omega_G$ with the property that $\int_G\omega_G=1$.
\end{proposition}
His proof starts: ``If $\dim G=0$, we just let $\omega_G$ be the constant
function $1/k$, where $k$ is the cardinality of $G$'' (p.~410); that is
Eq.~\eqref{haar:finite-average}. Otherwise it is the frame argument above,
followed by $\omega_G=(\int_G\omega)^{-1}\omega$. He then remarks (p.~411):
\begin{quote}
The orientation form whose existence is asserted in this proposition is
called the \emph{Haar volume form on $G$}. Similarly, the map
$f\mapsto\int_Gf\omega_G$ is called the \emph{Haar integral}. Observe that
the proof above did not use the fact that $G$ was compact until the last
paragraph; thus every Lie group has a left-invariant orientation form that
is uniquely defined up to a constant multiple. It is only in the compact
case, however, that we can use the volume normalization to single out a
unique one.
\end{quote}

\subsection{Right invariance: the modular function}\label{haar:sec:modular}
Left and right translations commute, $L_a\circ R_g=R_g\circ L_a$, so
$R_g^*\omega$ is again left invariant:
\begin{equation}
 L_a^*(R_g^*\omega)=(R_g\circ L_a)^*\omega=(L_a\circ R_g)^*\omega
 =R_g^*(L_a^*\omega)=R_g^*\omega .
\end{equation}
Hence, by the uniqueness up to a constant,
\begin{equation}\label{haar:modular-def}
 R_g^*\omega=\Delta(g)\,\omega,\qquad \Delta(g)\in\mathbb R\setminus\{0\}.
\end{equation}
\emph{Compactness forces $|\Delta(g)|=1$.} Measure volume by the density
$|\omega|$, which $R_g$ multiplies by $|\Delta(g)|$. The diffeomorphism $R_g$
maps $G$ onto $G$, so the total volume is unchanged, and the total volume is
finite and positive because $G$ is compact:
\begin{equation}\label{haar:modular-one}
 \int_G|\omega|=\int_GR_g^*|\omega|=|\Delta(g)|\int_G|\omega|
 \quad\Longrightarrow\quad |\Delta(g)|=1 .
\end{equation}
This is Kosmann-Schwarzbach's proof of Theorem~\ref{haar:thm:KS-2-1}(iii),
where $\mu_h(1)=\mu(1)=\Delta(h)\mu(1)$ (p.~35), and Hall's argument in
Appendix~D (p.~435), quoted with $\alpha$ the normalized left-invariant form:
\begin{quote}
If we translate $\alpha$ on the right by some $x\in K$, the resulting form
$\alpha^x$ is easily seen to be, again, a left-invariant volume form, which
must agree with $\alpha$ up to a constant. On the other hand, $\alpha^x$ is
still normalized, so it must actually agree with $\alpha$. Similarly, the
pullback of $\alpha$ by the map $x\mapsto x^{-1}$ is easily seen to be
left-invariant and normalized and thus coincides with $\alpha$. Thus,
$\alpha$ is invariant under both left and right translations and under
inversions.
\end{quote}
Hall's $K$ is connected. Then $\Delta$ is a continuous function with values
in $\{\pm1\}$ and $\Delta(e)=1$, so $\Delta\equiv1$ and the form itself is
bi-invariant.

The constant can be computed. Evaluate Eq.~\eqref{haar:modular-def} at $e$,
using Eq.~\eqref{haar:omega-def} at $V=g$ and
$(L_{g^{-1}}\circ R_g)(V)=g^{-1}Vg$, whose derivative at $e$ is
$\operatorname{Ad}_{g^{-1}}$:
\begin{align}
 (R_g^*\omega)_e&=(dR_g)_e^*\,\omega_g
 =(dR_g)_e^*\,(dL_{g^{-1}})_g^*\,\omega_e
 =\big(d(L_{g^{-1}}\circ R_g)\big)_e^*\,\omega_e\notag\\
 &=\operatorname{Ad}_{g^{-1}}^*\,\omega_e
 =\det(\operatorname{Ad}_{g^{-1}})\,\omega_e .\label{haar:modular}
\end{align}
The last step is the action of a linear map on a top-degree form. Thus
$\Delta(g)=\det\operatorname{Ad}_{g^{-1}}$, and since $\operatorname{Ad}$ is
a homomorphism and $\det$ is multiplicative, $\Delta$ is a homomorphism,
$\Delta(gg')=\Delta(g)\Delta(g')$. Kosmann-Schwarzbach calls it the
\emph{modular function} and the group \emph{unimodular} when $\Delta=1$
(p.~35).

The sign of $\Delta(g)$ can be $-1$ on a disconnected group. In $O(2)$, a
reflection $s$ acts on the one-dimensional $\mathfrak{so}(2)$ by
$\operatorname{Ad}_sX=sXs^{-1}=-X$, so $\Delta(s)=-1$ and $R_s$ reverses the
orientation fixed by $\omega$. For this reason the measure is defined by the
density $|\omega|$. Change variables by $R_g$:
\begin{equation}\label{haar:right-invariant-integral}
 \int_G(f\circ R_g)\,|\omega|
 =\frac1{|\Delta(g)|}\int_G(f\circ R_g)\,R_g^*|\omega|
 =\frac1{|\Delta(g)|}\int_Gf\,|\omega|
 =\int_Gf\,|\omega| .
\end{equation}
Equation~\eqref{haar:left-invariant-integral} holds with $|\omega|$ in place
of $\omega$ by the same steps, and
\begin{equation}\label{haar:normalized}
 \int_Gf\,dh:=\frac{\int_Gf\,|\omega|}{\int_G|\omega|}
\end{equation}
satisfies Eq.~\eqref{haar:bi-invariance}. Without compactness
Eq.~\eqref{haar:modular-one} is unavailable and $\Delta$ can differ from
$1$; the affine group of Exercise~3.7(b) below is an example.

\subsection{Uniqueness and inversion}
Let $h$ be the bi-invariant probability measure of
Eq.~\eqref{haar:normalized}, and let $\mu$ be any Borel probability measure
that is only left invariant. For continuous $f$, the function
$(V,W)\mapsto f(VW)$ is continuous on the compact space $G\times G$, so the
iterated integrals of it against $h\otimes\mu$ agree (Fubini). Integrate
first over $V$ with right invariance of $h$, then over $W$ with
left invariance of $\mu$:
\begin{align}
 \int_G\Big(\int_Gf(VW)\,h(dV)\Big)\mu(dW)
 &=\int_G\Big(\int_Gf(V)\,h(dV)\Big)\mu(dW)=\int_Gf\,dh,\\
 \int_G\Big(\int_Gf(VW)\,\mu(dW)\Big)h(dV)
 &=\int_G\Big(\int_Gf(W)\,\mu(dW)\Big)h(dV)=\int_Gf\,d\mu .
\end{align}
The left sides are equal, so $\int f\,d\mu=\int f\,dh$ for every continuous
$f$, and $\mu=h$. A left-invariant normalized measure is automatically the
bi-invariant one. This proves Theorem~\ref{haar:thm:KS-2-1}(ii) and (iii)
for compact Lie groups.

For inversion define $\mu(f)=\int_Gf(V^{-1})\,h(dV)$. Right invariance of
$h$, applied with $g_R=g^{-1}$, shows that $\mu$ is left invariant:
\begin{equation}\label{haar:inversion}
 \mu(f\circ L_g)=\int_Gf(gV^{-1})\,h(dV)
 =\int_Gf\big((Vg^{-1})^{-1}\big)\,h(dV)
 =\int_Gf(V^{-1})\,h(dV)=\mu(f).
\end{equation}
By uniqueness $\mu=h$. This completes Theorem~\ref{haar:thm:exist-unique}.
Uniqueness is what makes the examples below cheap: any left-invariant
probability measure that one recognizes is \emph{the} Haar measure.

\subsection{From an integral to a measure}\label{haar:sec:riesz}
Hall continues in Appendix~D (pp.~435--436):
\begin{quote}
Now, integration of a smooth function $f$ against $\alpha$ satisfies
\[
 \Big|\int_Kf\alpha\Big|\le\sup_K|f|.
\]
Meanwhile, by the Stone--Weierstrass theorem (Theorem~7.33 in [Rud1]), every
continuous function on $K$ can be uniformly approximated by smooth
functions. Thus, the map $f\mapsto\int_Kf\alpha$ extends by continuity from
smooth functions to continuous functions, and if $f$ is non-negative,
$\int_Kf\alpha$ will be non-negative. It then follows from the Riesz
representation theorem that there is a unique measure $\mu$ on the Borel
$\sigma$-algebra in $K$ such that
\[
 \int_Kf\alpha=\int_Kf(x)\,d\mu(x)
\]
for all continuous functions $f$ on $K$. (See Theorems~2.14 and 2.18 in
[Rud2]).
\end{quote}
The inequality is positivity and normalization. Since $\alpha$ is a
positively oriented form, $\int_Kg\alpha\ge0$ for $g\ge0$ (Lee,
Prop.~16.6(c)). Apply this to $g=\sup|f|-|f|\ge0$ and use $\int_K\alpha=1$:
\begin{equation}
 \Big|\int_Kf\alpha\Big|\le\int_K|f|\,\alpha\le\sup_K|f|\int_K\alpha=\sup_K|f| .
\end{equation}
The Riesz theorem converts the positive linear functional
$f\mapsto\int_Kf\alpha$ with value $1$ at $f=1$ into the Borel probability
measure of Definition~\ref{haar:def:haar}.

\subsection{Examples: \texorpdfstring{$\mathbb Z_2$}{Z2}, \texorpdfstring{$U(1)$}{U(1)}, \texorpdfstring{$O(2)$}{O(2)}, the affine group, \texorpdfstring{$SU(2)$}{SU(2)}}
\paragraph{$\mathbb Z_2=\{+1,-1\}$.} Here $|G|=2$, and
Eq.~\eqref{haar:finite-average} gives $h(\{+1\})=h(\{-1\})=\tfrac12$. For
$N$ independent links the product is $2^{-N}$ times counting measure.

\paragraph{$U(1)$.} Write $V=e^{i\theta}$, $0\le\theta<2\pi$. Left
multiplication by $e^{i\alpha}$ is $\theta\mapsto\theta+\alpha$ modulo
$2\pi$, which preserves Lebesgue measure on the circle. Kosmann-Schwarzbach
writes the result in two ways (p.~35). With $z=e^{i\theta}$,
$dz=ie^{i\theta}\,d\theta$, so
\begin{equation}\label{haar:U1}
 h(dV)=\frac{d\theta}{2\pi}=\frac{ie^{i\theta}\,d\theta}{2\pi i\,e^{i\theta}}
 =\frac{dz}{2i\pi z}.
\end{equation}

\begin{exercise}[Kosmann-Schwarzbach, Exercise~3.7, p.~45 \cite{dgdt:YKosmann-Schwarzbach2010}]
\emph{Haar measure.}
(a) Find the Haar measure on $O(2)$.
(b) Find the left-invariant and the right-invariant measures on the group of
affine transformations of $\mathbb R$.
\end{exercise}
\begin{proof}[Solution]
(a) Let $R_\theta$ be the rotation through $\theta$ and
$s=\operatorname{diag}(1,-1)$. Every element of $O(2)$ is $R_\theta$ or
$sR_\theta$, and $sR_\alpha s^{-1}=R_{-\alpha}$. The candidate is
\begin{equation}\label{haar:O2}
 \int_{O(2)}f\,dh=\frac12\int_0^{2\pi}\big(f(R_\theta)+f(sR_\theta)\big)
 \frac{d\theta}{2\pi}.
\end{equation}
Left multiplication by $R_\alpha$ sends $R_\theta\mapsto R_{\theta+\alpha}$
and $sR_\theta\mapsto R_\alpha sR_\theta=s\,(s^{-1}R_\alpha s)R_\theta
=sR_{\theta-\alpha}$. Both are translations of $\theta$ modulo $2\pi$ and
preserve $d\theta$. Left multiplication by $s$ sends $R_\theta\mapsto
sR_\theta$ and $sR_\theta\mapsto R_\theta$, exchanging the two terms. These
generate all left multiplications, so Eq.~\eqref{haar:O2} is a left-invariant
probability measure, hence the Haar measure. The reflection $s$ reverses the
orientation of $\omega$ (Section~\ref{haar:sec:modular}); the measure is
unaffected because it is built from $|\omega|$.

(b) Write $(a,b)$, $a\ne0$, for $x\mapsto ax+b$. Composition gives
$(a,b)(c,d)\colon x\mapsto a(cx+d)+b$, that is
\begin{equation}
 (a,b)(c,d)=(ac,\ ad+b).
\end{equation}
\emph{Left.} Left multiplication by $(c,d)$ is
$(a,b)\mapsto(ca,\ cb+d)$, with Jacobian determinant $c^2$. Substitute
$a'=ca$, $b'=cb+d$, so $da'\,db'=c^2\,da\,db$ and $a'^2=c^2a^2$:
\begin{equation}
 \int f(ca,cb+d)\,\frac{da\,db}{a^2}
 =\int f(a',b')\,\frac{c^{-2}\,da'\,db'}{c^{-2}a'^2}
 =\int f(a',b')\,\frac{da'\,db'}{a'^2}.
\end{equation}
\emph{Right.} Right multiplication by $(c,d)$ is
$(a,b)\mapsto(ac,\ ad+b)$, with Jacobian matrix
$\begin{bmatrix}c&0\\d&1\end{bmatrix}$ and determinant $c$. Substitute
$a'=ac$, $b'=ad+b$, so $da'\,db'=|c|\,da\,db$ and $|a'|=|c||a|$:
\begin{equation}
 \int f(ac,ad+b)\,\frac{da\,db}{|a|}
 =\int f(a',b')\,\frac{|c|^{-1}\,da'\,db'}{|c|^{-1}|a'|}
 =\int f(a',b')\,\frac{da'\,db'}{|a'|}.
\end{equation}
So $\frac{da\,db}{a^2}$ is left invariant and $\frac{da\,db}{|a|}$ is right
invariant. They are not multiples of each other. In
Kosmann-Schwarzbach's convention $\mu(f\circ r_h)=\Delta(h)\mu(f)$, right
multiplication applied to the left measure gives, with the same substitution
$a'=ac$, $b'=ad+b$,
\begin{equation}
 \int f(ac,ad+b)\,\frac{da\,db}{a^2}
 =\int f(a',b')\,\frac{|c|^{-1}c^2}{a'^2}\,da'\,db'
 =|c|\int f(a',b')\,\frac{da'\,db'}{a'^2},
\end{equation}
so $\Delta(c,d)=|c|$. The group is not unimodular.
\end{proof}

\paragraph{$SU(2)$.} By Eq.~\eqref{quat:su2-general}, $SU(2)$ is the set
of matrices of norm one in the real four-dimensional space
\begin{equation}\label{haar:quaternion-space}
 \mathcal M=\Big\{X=\begin{bmatrix}\alpha&-\bar\beta\\ \beta&\bar\alpha\end{bmatrix}
 :\alpha,\beta\in\mathbb C\Big\},\qquad
 \det X=|\alpha|^2+|\beta|^2=:|X|^2 ,
\end{equation}
and $|X|$ is the Euclidean norm of the real and imaginary parts of $\alpha$
and $\beta$ in $\mathbb R^4$. The space $\mathcal M$ is closed under
multiplication:
\begin{equation}\label{haar:M-product}
 \begin{bmatrix}\alpha&-\bar\beta\\ \beta&\bar\alpha\end{bmatrix}
 \begin{bmatrix}\gamma&-\bar\delta\\ \delta&\bar\gamma\end{bmatrix}
 =\begin{bmatrix}\alpha\gamma-\bar\beta\delta&-(\overline{\beta\gamma+\bar\alpha\delta})\\
 \beta\gamma+\bar\alpha\delta&\overline{\alpha\gamma-\bar\beta\delta}\end{bmatrix}.
\end{equation}
In the upper right, $-\alpha\bar\delta-\bar\beta\bar\gamma
=-(\overline{\bar\alpha\delta+\beta\gamma})$; in the lower right,
$-\beta\bar\delta+\bar\alpha\bar\gamma=\overline{\alpha\gamma-\bar\beta\delta}$.
For fixed $U\in SU(2)$ the maps $X\mapsto UX$ and $X\mapsto XU$ are
therefore real-linear maps of $\mathcal M\cong\mathbb R^4$, and
$|UX|^2=\det U\det X=|X|^2$. They are orthogonal maps of $\mathbb R^4$
and restrict to isometries of the unit sphere $S^3$. An isometry preserves
the Riemannian volume measure $\sigma$ of $S^3$, so $\sigma$ is invariant
under left and right multiplication. By uniqueness,
\begin{equation}\label{haar:SU2-round}
 h=\frac{\sigma}{\sigma(S^3)}=\frac{\sigma}{2\pi^2}
 \qquad\text{on }SU(2)\approx S^3 .
\end{equation}
The total volume $2\pi^2$ is computed in the next exercise.

\begin{exercise}[Kosmann-Schwarzbach, Exercise~5.4, p.~79 \cite{dgdt:YKosmann-Schwarzbach2010}]
\emph{Haar measure on the group $SU(2)$.}
We parametrize $SU(2)$ off a set of measure zero by setting, for an element
$g\in SU(2)$, $g=\begin{pmatrix}a&b\\-\bar b&\bar a\end{pmatrix}$,
$a=a_1+ia_2$ and $b=b_1+ib_2$, and
\[
 a_1=\cos\theta_3,\quad a_2=\sin\theta_3\cos\theta_2,\quad
 b_1=\sin\theta_3\sin\theta_2\cos\theta_1,\quad
 b_2=\sin\theta_3\sin\theta_2\sin\theta_1,
\]
for $\theta_1\in(0,2\pi)$, $\theta_2\in(0,\pi)$, $\theta_3\in(0,\pi)$.

(a) Show that the Haar integral on $SU(2)$ can be written
\[
 I(f)=\frac{1}{2\pi^2}\int_{\mathcal O}f(\theta_1,\theta_2,\theta_3)
 \sin^2\theta_3\sin\theta_2\,d\theta_1\,d\theta_2\,d\theta_3,
\]
where $\mathcal O=(0,2\pi)\times(0,\pi)\times(0,\pi)$.

(b) Set $\theta_3=t$. Show that the Haar integral of a class function $f$ is
\[
 \frac{2}{\pi}\int_0^\pi f(t)\sin^2t\,dt.
\]
\end{exercise}
\begin{proof}[Solution]
Kosmann-Schwarzbach's $g$ is our matrix of Eq.~\eqref{haar:quaternion-space}
with $\alpha=a$ and $\beta=-\bar b$; the map $(a_1,a_2,b_1,b_2)\mapsto
(\operatorname{Re}\alpha,\operatorname{Im}\alpha,\operatorname{Re}\beta,
\operatorname{Im}\beta)=(a_1,a_2,-b_1,b_2)$ is orthogonal, so
Eq.~\eqref{haar:SU2-round} holds in the coordinates
$x=(a_1,a_2,b_1,b_2)\in S^3\subset\mathbb R^4$.

(a) Write $x=(\cos\theta_3,\ \sin\theta_3\,\mathbf n)$ with
$\mathbf n=(\cos\theta_2,\ \sin\theta_2\cos\theta_1,\ \sin\theta_2\sin\theta_1)\in S^2$.
Differentiate, using $\mathbf n\cdot\mathbf n=1$:
\begin{align}
 dx&=\big(-\sin\theta_3\,d\theta_3,\ \cos\theta_3\,\mathbf n\,d\theta_3
     +\sin\theta_3\,d\mathbf n\big),\qquad
 \mathbf n\cdot d\mathbf n=\tfrac12d(\mathbf n\cdot\mathbf n)=0,\notag\\
 dx\cdot dx&=\sin^2\theta_3\,d\theta_3^2+\cos^2\theta_3\,d\theta_3^2
     +2\sin\theta_3\cos\theta_3\,d\theta_3\,(\mathbf n\cdot d\mathbf n)
     +\sin^2\theta_3\,d\mathbf n\cdot d\mathbf n\notag\\
 &=d\theta_3^2+\sin^2\theta_3\,d\mathbf n\cdot d\mathbf n .
 \label{haar:S3-metric-1}
\end{align}
For the two-sphere,
$\frac{\partial\mathbf n}{\partial\theta_2}=(-\sin\theta_2,\
\cos\theta_2\cos\theta_1,\ \cos\theta_2\sin\theta_1)$ and
$\frac{\partial\mathbf n}{\partial\theta_1}=(0,\ -\sin\theta_2\sin\theta_1,\
\sin\theta_2\cos\theta_1)$ have squared lengths $1$ and $\sin^2\theta_2$, and
their dot product is
$-\cos\theta_2\sin\theta_2\cos\theta_1\sin\theta_1
+\cos\theta_2\sin\theta_2\sin\theta_1\cos\theta_1=0$. So
\begin{equation}\label{haar:S3-metric}
 dx\cdot dx=d\theta_3^2+\sin^2\theta_3\big(d\theta_2^2
 +\sin^2\theta_2\,d\theta_1^2\big).
\end{equation}
The metric is diagonal, and the volume density is the product of the square
roots of the diagonal entries:
\begin{equation}\label{haar:S3-volume}
 \sigma(dx)=\sin^2\theta_3\sin\theta_2\,d\theta_1\,d\theta_2\,d\theta_3,
 \qquad
 \sigma(S^3)=\int_0^{2\pi}d\theta_1\int_0^\pi\sin\theta_2\,d\theta_2
 \int_0^\pi\sin^2\theta_3\,d\theta_3=2\pi\cdot2\cdot\frac\pi2=2\pi^2 .
\end{equation}
The last factor uses $\sin^2\theta_3=\frac12(1-\cos2\theta_3)$. The excluded
coordinate boundaries have measure zero. Dividing by $2\pi^2$ as in
Eq.~\eqref{haar:SU2-round} gives $I(f)$.

(b) Every $g\in SU(2)$ is unitary, so $g=W\operatorname{diag}(\lambda_1,
\lambda_2)W^{-1}$ with $W\in U(2)$ and $|\lambda_i|=1$. Replacing $W$ by
$W/\sqrt{\det W}$ (either root) gives $W\in SU(2)$ and does not change the
product, since a scalar commutes with the diagonal matrix. With
$\lambda_1\lambda_2=\det g=1$ and $\operatorname{tr}g=a+\bar a=2a_1$ real,
\begin{equation}\label{haar:SU2-classes}
 \lambda_{1,2}=e^{\pm it},\qquad \operatorname{tr}g=2\cos t,\qquad
 t=\theta_3\in[0,\pi].
\end{equation}
A class function is constant on conjugacy classes, hence a function of $t$
alone. Integrate out $\theta_1$ and $\theta_2$ in (a):
\begin{equation}\label{haar:weyl-SU2}
 \int_{SU(2)}f\,dh=\frac{1}{2\pi^2}\cdot2\pi\cdot2\int_0^\pi f(t)\sin^2t\,dt
 =\frac2\pi\int_0^\pi f(t)\sin^2t\,dt .
\end{equation}
\end{proof}
The weight $\sin^2t$ vanishes at $t=0$ and $t=\pi$: the conjugacy class at
fixed $t$ is the two-sphere $\{(\cos t,\sin t\,\mathbf n)\}$ of radius
$\sin t$ in Eq.~\eqref{haar:S3-metric}, which shrinks to the points $\pm1$.

\subsection{Euler angles on \texorpdfstring{$SO(3)$}{SO(3)} and \texorpdfstring{$SU(2)$}{SU(2)}}\label{haar:sec:euler}
Kosmann-Schwarzbach's $\eta_k$ are the matrices $E_k$ of
Eq.~\eqref{quat:Ei-entries}, $\eta_k=\widehat{\mathbf e_k}$, so
$\exp(t\eta_k)$ is the rotation through $t$ about $\mathbf e_k$ in
Eq.~\eqref{quat:rotate-about-n}. Her $\xi_k$ are those of
Eq.~\eqref{quat:KS-xi}, and $\varphi\colon SU(2)\to SO(3)$ is the covering
map of Eq.~\eqref{quat:KS-phi}, with $d\varphi(\xi_k)=\eta_k$.

\begin{exercise}[Kosmann-Schwarzbach, Exercise~5.5, p.~80 \cite{dgdt:YKosmann-Schwarzbach2010}]
\emph{Euler angles and Haar measure on the group $SO(3)$.}
(a) Show that each rotation $g\in SO(3)$ can be written
\[
 g=\exp(\alpha\eta_3)\exp(\beta\eta_1)\exp(\gamma\eta_3)
\]
with $0\le\alpha<2\pi$, $0\le\beta\le\pi$, $0\le\gamma<2\pi$, uniquely if
$\beta\ne0$ and $\beta\ne\pi$. The angles $\alpha$, $\beta$, $\gamma$ are
called the \emph{Euler angles} of the rotation.
(b) Show that the Haar integral of $SO(3)$ can be written
\[
 I(f)=\frac1{8\pi^2}\int_{\Omega_0}f(\alpha,\beta,\gamma)\sin\beta\,
 d\alpha\,d\beta\,d\gamma,
\]
where $\Omega_0=(0,2\pi)\times(0,\pi)\times(0,2\pi)$.
\end{exercise}
\begin{proof}[Solution]
(a) By Eq.~\eqref{quat:rotate-about-n}, $\exp(\beta\eta_1)$ sends
$\mathbf e_3\mapsto\cos\beta\,\mathbf e_3+\sin\beta\,\mathbf e_1\times\mathbf e_3
=\cos\beta\,\mathbf e_3-\sin\beta\,\mathbf e_2$, and $\exp(\alpha\eta_3)$
sends $\mathbf e_2\mapsto\cos\alpha\,\mathbf e_2-\sin\alpha\,\mathbf e_1$
and fixes $\mathbf e_3$. Since $\exp(\gamma\eta_3)\mathbf e_3=\mathbf e_3$,
\begin{equation}\label{haar:euler-axis}
 \exp(\alpha\eta_3)\exp(\beta\eta_1)\exp(\gamma\eta_3)\,\mathbf e_3
 =\cos\beta\,\mathbf e_3-\sin\beta\,(\cos\alpha\,\mathbf e_2-\sin\alpha\,\mathbf e_1)
 =(\sin\beta\sin\alpha,\ -\sin\beta\cos\alpha,\ \cos\beta).
\end{equation}
Given $g$, let $\mathbf u=g\mathbf e_3$, a unit vector. Take
$\beta=\arccos u^3\in[0,\pi]$. If $0<\beta<\pi$, then $\sin\beta>0$ and
$(\sin\alpha,-\cos\alpha)=(u^1,u^2)/\sin\beta$ is a unit vector, which
determines $\alpha\in[0,2\pi)$ uniquely. Put
$k=\big(\exp(\alpha\eta_3)\exp(\beta\eta_1)\big)^{-1}g$. By
Eq.~\eqref{haar:euler-axis} with $\gamma=0$, $k\mathbf e_3=\mathbf e_3$. So
$k\in SO(3)$ maps the plane $\mathbf e_3^\perp$ to itself with determinant
$1$, and is a rotation $\exp(\gamma\eta_3)$ for a unique $\gamma\in[0,2\pi)$.
Conversely, any factorization has $g\mathbf e_3$ given by
Eq.~\eqref{haar:euler-axis}, so for $0<\beta<\pi$ it fixes $\beta$, then
$\alpha$, then $\gamma$. If $\beta=0$ or $\pi$ the vector
$g\mathbf e_3=\pm\mathbf e_3$ does not determine $\alpha$, and only
$\alpha\pm\gamma$ is determined.

(b) Write $g=ABC$ with $A=\exp(\alpha\eta_3)$, $B=\exp(\beta\eta_1)$,
$C=\exp(\gamma\eta_3)$. Then $A^{-1}dA=\eta_3\,d\alpha$,
$B^{-1}dB=\eta_1\,d\beta$, $C^{-1}dC=\eta_3\,d\gamma$, and by the product
rule
\begin{align}
 \Theta=g^{-1}dg
 &=C^{-1}B^{-1}A^{-1}\big(dA\,BC+A\,dB\,C+AB\,dC\big)\notag\\
 &=C^{-1}B^{-1}(A^{-1}dA)BC+C^{-1}(B^{-1}dB)C+C^{-1}dC\notag\\
 &=\operatorname{Ad}_{C^{-1}}\big(\operatorname{Ad}_{B^{-1}}\eta_3\,d\alpha
   +\eta_1\,d\beta+\eta_3\,d\gamma\big),\label{haar:euler-MC}
\end{align}
where $\operatorname{Ad}_{C^{-1}}\eta_3=\eta_3$ because $C$ commutes with
$\eta_3$. For $R\in SO(3)$, $\operatorname{Ad}_R\eta_k=R\,\widehat{\mathbf e_k}
\,R^{-1}=\widehat{R\mathbf e_k}$, so on coefficient vectors in the basis
$\eta_k$, $\operatorname{Ad}_R$ acts as the matrix $R$. Here $B^{-1}$ is the
rotation through $-\beta$ about $\mathbf e_1$, and by
Eq.~\eqref{quat:rotate-about-n}
\begin{equation}
 B^{-1}\mathbf e_3=\cos\beta\,\mathbf e_3-\sin\beta\,\mathbf e_1\times\mathbf e_3
 =\sin\beta\,\mathbf e_2+\cos\beta\,\mathbf e_3,
 \qquad \operatorname{Ad}_{B^{-1}}\eta_3=\sin\beta\,\eta_2+\cos\beta\,\eta_3 .
\end{equation}
Write $\Theta=\Theta^k\eta_k$ and $\Theta^k=(C^{-1})^k{}_lM^l{}_j\,dq^j$ with
$(q^1,q^2,q^3)=(\alpha,\beta,\gamma)$. The columns of $M$ are the
coefficient vectors of $\operatorname{Ad}_{B^{-1}}\eta_3$, $\eta_1$ and
$\eta_3$ in Eq.~\eqref{haar:euler-MC}:
\begin{equation}\label{haar:euler-M}
 M=\begin{bmatrix}0&1&0\\ \sin\beta&0&0\\ \cos\beta&0&1\end{bmatrix},
 \qquad \det M=-1\cdot\det\begin{bmatrix}\sin\beta&0\\ \cos\beta&1\end{bmatrix}
 =-\sin\beta ,
\end{equation}
expanding along the first row. By Eq.~\eqref{haar:maurer-cartan} and
$\det C^{-1}=1$,
\begin{equation}\label{haar:euler-volume}
 \omega=\Theta^1\wedge\Theta^2\wedge\Theta^3
 =\det(C^{-1})\det M\,d\alpha\wedge d\beta\wedge d\gamma
 =-\sin\beta\,d\alpha\wedge d\beta\wedge d\gamma .
\end{equation}
By (a) the map $(\alpha,\beta,\gamma)\mapsto g$ is a diffeomorphism from
$\Omega_0$ onto $SO(3)$ minus a set of measure zero, so $|\omega|$ integrates
as $\int_{\Omega_0}f\sin\beta\,d\alpha\,d\beta\,d\gamma$ (Lee,
Prop.~16.8). The total is
$\int_0^{2\pi}d\alpha\int_0^\pi\sin\beta\,d\beta\int_0^{2\pi}d\gamma
=2\pi\cdot2\cdot2\pi=8\pi^2$, and Eq.~\eqref{haar:normalized} gives $I(f)$.
\end{proof}

\begin{exercise}[Kosmann-Schwarzbach, Exercise~5.6, p.~80 \cite{dgdt:YKosmann-Schwarzbach2010}]
\emph{Euler angles and Haar measure on the group $SU(2)$.}
Let $\Omega=(0,2\pi)\times(0,\pi)\times(-2\pi,2\pi)$.
(a) Show that for every pair of complex numbers $a$ and $b$ such that
$|a|^2+|b|^2=1$, $\mathcal Im(a)\ne0$, and $\mathcal Re(b)\ne0$, there is a
unique triple $(\phi,\theta,\psi)\in\Omega$ such that
$a=\cos\frac\theta2e^{i\phi+\psi/2}$, $b=i\sin\frac\theta2e^{i\phi-\psi/2}$.
Conclude that off a set of measure zero of $SU(2)$, each matrix
$g=\left(\begin{smallmatrix}a&b\\-\bar b&\bar a\end{smallmatrix}\right)\in SU(2)$
can be written in one and only one way as
\[
 g=\exp(\phi\xi_3)\exp(\theta\xi_1)\exp(\psi\xi_3),
\]
with $(\phi,\theta,\psi)\in\Omega$. Show that the image of
$g(\phi,\theta,\psi)$ under the morphism $\varphi$ of $SO(3)$ onto $SU(2)$
is the rotation with Euler angles $(\phi,\theta,\psi)$.
(b) Show that the Haar integral on $SU(2)$ can be written
\[
 I(f)=\frac1{16\pi^2}\int_\Omega f(\phi,\theta,\psi)\sin\theta\,
 d\phi\,d\theta\,d\psi.
\]
\end{exercise}
\emph{Editorial corrections.} The printed exponents $i\phi\pm\psi/2$ should
be $i(\phi\pm\psi)/2$, as the product below shows. Her $\varphi$ maps
$SU(2)$ onto $SO(3)$ (Eq.~\eqref{quat:KS-phi}). The two conditions on $a$
and $b$ do not exclude the boundary $\phi=0$ (for example
$\arg a=\frac\pi4$, $\arg(b/i)=-\frac\pi4$); the solution states the
exceptional set that the argument requires.
\begin{proof}[Solution]
(a) With $\xi_3=\frac i2\sigma_3$ and $\xi_1=\frac i2\sigma_1$ from
Eq.~\eqref{quat:KS-xi}, and $\sigma_1^2=\mathbf 1$,
$\exp(\phi\xi_3)=\operatorname{diag}(e^{i\phi/2},e^{-i\phi/2})$ and
$\exp(\theta\xi_1)=\cos\frac\theta2\,\mathbf 1+i\sin\frac\theta2\,\sigma_1$.
Write $c=\cos\frac\theta2$, $s=\sin\frac\theta2$. Multiplying on the left by
the first diagonal matrix scales the rows, and on the right by the second
scales the columns:
\begin{align}
 \exp(\phi\xi_3)\exp(\theta\xi_1)\exp(\psi\xi_3)
 &=\begin{bmatrix}ce^{i\phi/2}&ise^{i\phi/2}\\ ise^{-i\phi/2}&ce^{-i\phi/2}\end{bmatrix}
   \begin{bmatrix}e^{i\psi/2}&0\\0&e^{-i\psi/2}\end{bmatrix}\notag\\
 &=\begin{bmatrix}ce^{i(\phi+\psi)/2}&ise^{i(\phi-\psi)/2}\\
   ise^{-i(\phi-\psi)/2}&ce^{-i(\phi+\psi)/2}\end{bmatrix}.\label{haar:SU2-euler}
\end{align}
This has the form $\left(\begin{smallmatrix}a&b\\-\bar b&\bar a\end{smallmatrix}\right)$
with
\begin{equation}\label{haar:SU2-euler-ab}
 a=\cos\tfrac\theta2\,e^{i(\phi+\psi)/2},\qquad
 b=i\sin\tfrac\theta2\,e^{i(\phi-\psi)/2},\qquad
 -\bar b=is\,e^{-i(\phi-\psi)/2}.
\end{equation}
Conversely let $ab\ne0$. Then $\cos\frac\theta2=|a|\in(0,1)$ fixes
$\theta\in(0,\pi)$, and $\sin\frac\theta2=|b|$. Let $A=\arg a$ and
$B=\arg(b/i)$, each defined modulo $2\pi$. We need
$\frac{\phi+\psi}2\equiv A$ and $\frac{\phi-\psi}2\equiv B$ modulo $2\pi$.
Adding, $\phi\equiv A+B$ modulo $2\pi$, which has one solution in
$(0,2\pi)$ unless $A+B\equiv0$. Given $\phi$, the first condition reads
$\psi\equiv2A-\phi$ modulo $4\pi$, which has one solution in $(-2\pi,2\pi)$
unless $\psi\equiv2\pi$ modulo $4\pi$. This pair satisfies the second
condition too: $\frac{\phi-\psi}2=\phi-\frac{\phi+\psi}2\equiv(A+B)-A=B$.
The excluded set $\{ab=0\}\cup\{\phi=0\}\cup\{\psi=\pm2\pi\}$ is a finite
union of lower-dimensional submanifolds, of measure zero.
Finally $\varphi$ is a homomorphism with $d\varphi(\xi_k)=\eta_k$, so
$\varphi(\exp(t\xi_k))=\exp(t\,d\varphi(\xi_k))=\exp(t\eta_k)$ and
\begin{equation}
 \varphi(g)=\exp(\phi\eta_3)\exp(\theta\eta_1)\exp(\psi\eta_3),
\end{equation}
the rotation with Euler angles $(\phi,\theta,\psi)$, with $\psi$ read modulo
$2\pi$.

(b) The computation of Eq.~\eqref{haar:euler-MC} applies with
$(\alpha,\beta,\gamma,\eta_k)$ replaced by $(\phi,\theta,\psi,\xi_k)$. By
the definition of $\varphi$ (\S\ref{quat:sec:KS-phi}), the matrix of
$\operatorname{Ad}_u$ in the basis $\xi_k$ is $\varphi(u)$. Here
$\varphi(B^{-1})=\exp(-\theta\eta_1)$, so
$\operatorname{Ad}_{B^{-1}}\xi_3=\sin\theta\,\xi_2+\cos\theta\,\xi_3$ and the
matrix $M$ is Eq.~\eqref{haar:euler-M} with $\beta\to\theta$. With
$\det\varphi(C^{-1})=1$,
\begin{equation}
 \Theta^1\wedge\Theta^2\wedge\Theta^3=-\sin\theta\,d\phi\wedge d\theta\wedge d\psi,
 \qquad
 \int_\Omega\sin\theta\,d\phi\,d\theta\,d\psi=2\pi\cdot2\cdot4\pi=16\pi^2 .
\end{equation}
By (a) the parametrization covers $SU(2)$ once off a set of measure zero,
and Eq.~\eqref{haar:normalized} gives $I(f)$.
\end{proof}
The two normalizations differ by the factor $2$ of the double cover. The
range of $\psi$ is $4\pi$ in $SU(2)$ and $2\pi$ in $SO(3)$.

\subsection{The Weyl integral formula}
Hall integrates class functions over a maximal torus $T$, a maximal subgroup
isomorphic to a product of circles (Definition~11.6, p.~312). The Weyl group
is $W=N(T)/T$, with $N(T)=\{x\in K:xTx^{-1}=T\}$ (Definition~11.8, p.~313).
Let $\mathfrak f$ be the orthogonal complement of $\mathfrak t$ in
$\mathfrak k$ and
\begin{equation}\label{haar:weyl-rho}
 \rho(t)=\det\big(\operatorname{Ad}'_{t^{-1}}-I\big),
\end{equation}
where $\operatorname{Ad}'_{t^{-1}}$ is the restriction of
$\operatorname{Ad}_{t^{-1}}$ to $\mathfrak f$ (Hall, Eq.~(11.12), p.~330).
\begin{theorem}[Hall, Theorem~11.30, Weyl Integral Formula, p.~330 \cite{BCHall2015}]\label{haar:thm:weyl}
For all continuous functions $f$ on $K$, we have
\[
 \int_Kf(x)\,dx=\frac1{|W|}\int_T\rho(t)\int_{K/T}f(yty^{-1})\,d[y]\,dt,
\]
where $dx$, $dt$, and $d[y]$ are the normalized, left-invariant volume forms
on $K$, $T$, and $K/T$, respectively and $|W|$ is the order of the Weyl group.
\end{theorem}
Hall's proof up to the constant computes the pullback of $dx$ under
$\Phi(t,[y])=yty^{-1}$, a map of degree $|W|$, and its Jacobian determinant is
$\rho(t)$. He fixes the constant in \S12.4. For a class function
(Definition~11.31, p.~331: $f(yxy^{-1})=f(x)$), $\int_{K/T}f(yty^{-1})\,d[y]=f(t)$
and the formula becomes his Corollary~11.32,
\begin{equation}\label{haar:weyl-class}
 \int_Kf(x)\,dx=\frac1{|W|}\int_T\rho(t)f(t)\,dt .
\end{equation}

\paragraph{$SU(2)$ (Hall, Example~11.33, p.~332).} Take $T$ the diagonal
subgroup, $t=\operatorname{diag}(e^{i\theta},e^{-i\theta})$. Under the
Hilbert--Schmidt inner product the complement of $\mathfrak t$ is
$\mathfrak f=\{X_z=\left(\begin{smallmatrix}0&z\\-\bar z&0\end{smallmatrix}\right):
z=x+iy\in\mathbb C\}$. Conjugate:
\begin{equation}
 \operatorname{Ad}_{t^{-1}}X_z=\begin{bmatrix}e^{-i\theta}&0\\0&e^{i\theta}\end{bmatrix}
 \begin{bmatrix}0&z\\-\bar z&0\end{bmatrix}
 \begin{bmatrix}e^{i\theta}&0\\0&e^{-i\theta}\end{bmatrix}
 =\begin{bmatrix}0&e^{-2i\theta}z\\-e^{2i\theta}\bar z&0\end{bmatrix}
 =X_{e^{-2i\theta}z}.
\end{equation}
So $\operatorname{Ad}'_{t^{-1}}$ is the rotation through $-2\theta$ of
$(x,y)\in\mathbb R^2$, and
\begin{equation}\label{haar:rho-SU2}
 \rho(t)=\det\begin{bmatrix}\cos2\theta-1&\sin2\theta\\-\sin2\theta&\cos2\theta-1\end{bmatrix}
 =(\cos2\theta-1)^2+\sin^22\theta=2-2\cos2\theta=4\sin^2\theta .
\end{equation}
The Weyl group has order $|W|=2$, and $dt=\frac{d\theta}{2\pi}$ by
Eq.~\eqref{haar:U1}. Equation~\eqref{haar:weyl-class} is Hall's
\begin{equation}\label{haar:hall-11-16}
 \int_{SU(2)}f(x)\,dx=\frac12\int_{-\pi}^{\pi}f\big(\operatorname{diag}(e^{i\theta},e^{-i\theta})\big)
 \,4\sin^2\theta\,\frac{d\theta}{2\pi}.
\end{equation}
Conjugation by $\left(\begin{smallmatrix}0&1\\-1&0\end{smallmatrix}\right)\in SU(2)$
exchanges $\operatorname{diag}(e^{i\theta},e^{-i\theta})$ and
$\operatorname{diag}(e^{-i\theta},e^{i\theta})$, so a class function is even
in $\theta$. With $\sin^2\theta$ also even,
\begin{equation}
 \frac12\int_{-\pi}^{\pi}f\,4\sin^2\theta\,\frac{d\theta}{2\pi}
 =\frac1\pi\int_{-\pi}^{\pi}f\,\sin^2\theta\,d\theta
 =\frac2\pi\int_0^\pi f\,\sin^2\theta\,d\theta ,
\end{equation}
which is Eq.~\eqref{haar:weyl-SU2}. The $S^3$ computation of Exercise~5.4
fixed the constant directly.

\subsection{Three integrals on \texorpdfstring{$SU(2)$}{SU(2)}}
\begin{proposition}\label{haar:prop:SU2-integrals}
For normalized Haar measure on $SU(2)$, every $X\in M_2(\mathbb C)$ and
every continuous class function $\phi$,
\begin{align}
 \int\operatorname{tr}V\,h(dV)&=0,\label{haar:trace-zero}\\
 \int(\operatorname{tr}V)^2\,h(dV)&=1,\label{haar:trace-square}\\
 \int VXV^{-1}\,h(dV)&=\frac{\operatorname{tr}X}{2}\,\mathbf 1,
 \label{haar:conjugation-average}\\
 \int\phi(W)\,W^{-1}\,h(dW)&=\Big(\frac12\int\phi(W)\operatorname{tr}W\,h(dW)\Big)\mathbf 1 .
 \label{haar:class-average}
\end{align}
\end{proposition}
\begin{proof}
Eq.~\eqref{haar:trace-zero}: $-\mathbf 1\in SU(2)$. Apply
Eq.~\eqref{haar:bi-invariance} with $g_L=-\mathbf 1$, $g_R=\mathbf 1$:
$\int\operatorname{tr}V\,h(dV)=\int\operatorname{tr}(-V)\,h(dV)
=-\int\operatorname{tr}V\,h(dV)$.

Eq.~\eqref{haar:trace-square}: insert $\operatorname{tr}V=2\cos t$ from
Eq.~\eqref{haar:SU2-classes} into Eq.~\eqref{haar:weyl-SU2}, and use
$4\cos^2t\sin^2t=\sin^22t=\frac12(1-\cos4t)$:
\begin{equation}
 \frac2\pi\int_0^\pi4\cos^2t\,\sin^2t\,dt
 =\frac2\pi\int_0^\pi\frac{1-\cos4t}{2}\,dt=\frac2\pi\cdot\frac\pi2=1 .
\end{equation}
Eq.~\eqref{haar:conjugation-average}: call the left side $A$. For
$U\in SU(2)$, substitute $V'=UV$, which is Eq.~\eqref{haar:bi-invariance}
with $g_L=U$:
\begin{equation}
 UAU^{-1}=\int(UV)X(UV)^{-1}\,h(dV)=\int V'X{V'}^{-1}\,h(dV')=A .
\end{equation}
Thus $A$ commutes with $\mathbf 1,i\sigma_1,i\sigma_2,i\sigma_3\in SU(2)$.
These span $M_2(\mathbb C)$ over $\mathbb C$, so $A$ commutes with every
$2\times2$ matrix and is a scalar, $A=c\mathbf 1$. The trace of the integrand
is $\operatorname{tr}(VXV^{-1})=\operatorname{tr}X$, so $2c=\operatorname{tr}X$.

Eq.~\eqref{haar:class-average}: call the left side $B$. For $U\in SU(2)$
substitute $W'=UWU^{-1}$ (Eq.~\eqref{haar:bi-invariance} with $g_L=U$,
$g_R=U^{-1}$), and use $\phi(U^{-1}W'U)=\phi(W')$:
\begin{equation}
 UBU^{-1}=\int\phi(W)\,(UWU^{-1})^{-1}\,h(dW)=\int\phi(W')\,{W'}^{-1}\,h(dW')=B .
\end{equation}
As before $B=c\mathbf 1$, and $2c=\int\phi(W)\operatorname{tr}W^{-1}\,h(dW)$.
For $W\in SU(2)$, $\operatorname{tr}W^{-1}=\operatorname{tr}W^\dagger
=\overline{\operatorname{tr}W}=\operatorname{tr}W$, since the trace is real by
Eq.~\eqref{haar:SU2-classes}.
\end{proof}
Equations~\eqref{haar:trace-zero} and \eqref{haar:trace-square} are the
statements $\langle\chi_0,\chi_{1/2}\rangle=0$ and
$\langle\chi_{1/2},\chi_{1/2}\rangle=1$ of Theorem~\ref{haar:thm:hall-12-15}
below. The argument for Eq.~\eqref{haar:conjugation-average} is Schur's
lemma (Theorem~\ref{haar:thm:schur}) for the defining representation, with
the spanning set written out.

\subsection{Averaging representations: unitarity, Schur, characters}
A finite-dimensional representation $\Pi$ of $K$ on $V$ is \emph{unitary}
for an inner product if each $\Pi(x)$ preserves it, and \emph{completely
reducible} if $V$ is a direct sum of irreducible invariant subspaces.

\begin{proposition}[Hall, Proposition~4.27, p.~92 \cite{BCHall2015}]\label{haar:prop:hall-4-27}
If $G$ is a matrix Lie group and $\Pi$ is a finite-dimensional unitary
representation of $G$, then $\Pi$ is completely reducible.
\end{proposition}
\begin{proof}
Let $W\subset V$ be invariant and $v\in W^\perp$. For $w\in W$,
$w'=\Pi(A^{-1})w\in W$, and unitarity gives
$\Pi(A)^*=\Pi(A)^{-1}=\Pi(A^{-1})$:
\begin{equation}
 \langle\Pi(A)v,w\rangle=\langle v,\Pi(A)^*w\rangle=\langle v,\Pi(A^{-1})w\rangle
 =\langle v,w'\rangle=0 .
\end{equation}
So $W^\perp$ is invariant and $V=W\oplus W^\perp$. Repeat on $W$ and
$W^\perp$; since $\dim V<\infty$ the process stops at irreducible summands.
\end{proof}
\begin{theorem}[Hall, Theorem~4.28, p.~92 \cite{BCHall2015}]\label{haar:thm:hall-4-28}
If $G$ is a compact matrix Lie group, every finite dimensional
representation of $G$ is completely reducible.
\end{theorem}
\begin{proof}
Hall's proof, ``Weyl's unitarian trick'' (pp.~93--94), uses a
right-invariant form $\alpha_A(Y_1,\dots,Y_k)=\alpha_I(Y_1A^{-1},\dots,Y_kA^{-1})$,
the mirror image of Eq.~\eqref{haar:omega-matrix}; by
Theorem~\ref{haar:thm:exist-unique} its normalization is $h$. Choose any
inner product on $V$ and average it:
\begin{equation}\label{haar:unitarian-trick}
 \langle v,w\rangle_G=\int_G\langle\Pi(A)v,\Pi(A)w\rangle\,h(dA).
\end{equation}
It is an inner product: the integrand of $\langle v,v\rangle_G$ is positive
and continuous for $v\ne0$. For $B\in G$, right invariance with $g_R=B$ gives
\begin{align}
 \langle\Pi(B)v,\Pi(B)w\rangle_G
 &=\int_G\langle\Pi(A)\Pi(B)v,\Pi(A)\Pi(B)w\rangle\,h(dA)
 =\int_G\langle\Pi(AB)v,\Pi(AB)w\rangle\,h(dA)\notag\\
 &=\int_G\langle\Pi(A)v,\Pi(A)w\rangle\,h(dA)=\langle v,w\rangle_G .
\end{align}
So $\Pi$ is unitary for $\langle\cdot,\cdot\rangle_G$, and
Proposition~\ref{haar:prop:hall-4-27} applies. Compactness makes the integral
finite.
\end{proof}
A linear map $\phi\colon V\to W$ between representations $(\Pi,V)$,
$(\Sigma,W)$ is \emph{intertwining} if $\phi\,\Pi(A)=\Sigma(A)\,\phi$ for all
$A$.
\begin{theorem}[Hall, Theorem~4.29, Schur's Lemma, p.~94 \cite{BCHall2015}]\label{haar:thm:schur}
\begin{enumerate}
\item Let $V$ and $W$ be irreducible real or complex representations of a
group or Lie algebra and let $\phi\colon V\to W$ be an intertwining map.
Then either $\phi=0$ or $\phi$ is an isomorphism.
\item Let $V$ be an irreducible complex representation of a group or Lie
algebra and let $\phi\colon V\to V$ be an intertwining map of $V$ with
itself. Then $\phi=\lambda I$, for some $\lambda\in\mathbb C$.
\end{enumerate}
\end{theorem}
\begin{proof}
(1) If $v\in\ker\phi$, then $\phi(\Pi(A)v)=\Sigma(A)\phi(v)=0$, so
$\ker\phi$ is invariant, hence $0$ or $V$. If $w=\phi(v)$, then
$\Sigma(A)w=\phi(\Pi(A)v)$, so the image is invariant, hence $0$ or $W$.
A nonzero $\phi$ is therefore injective and surjective.
(2) Over $\mathbb C$, $\phi$ has an eigenvalue $\lambda$ with eigenspace
$U\ne0$. For $u\in U$, $\phi\,\Pi(A)u=\Pi(A)\phi u=\lambda\Pi(A)u$, so $U$ is
invariant, $U=V$, and $\phi=\lambda I$.
\end{proof}
Consequently the intertwiners $V\to W$ form a space of dimension $1$ if
$V\cong W$ and $0$ otherwise. If $\phi_0\colon V\to W$ is an isomorphism and
$A$ any intertwiner, $\phi_0^{-1}A$ intertwines $V$ with itself, so
$\phi_0^{-1}A=\lambda I$ and $A=\lambda\phi_0$.

The \emph{character} of $(\Pi,V)$ is $\chi_\Pi(x)=\operatorname{tr}\Pi(x)$
(Hall, Definition~12.14, p.~351). It is a class function:
$\chi_\Pi(yxy^{-1})=\operatorname{tr}(\Pi(y)\Pi(x)\Pi(y)^{-1})
=\operatorname{tr}\Pi(x)$.
\begin{lemma}[Hall, Lemma~12.16, p.~352 \cite{BCHall2015}]\label{haar:lem:hall-12-16}
Suppose $(\Pi,V)$ is a finite-dimensional representation of $K$, and let
$P$ be the operator on $V$ given by
\[
 P=\int_K\Pi(x)\,dx.
\]
Then $P$ is a projection onto $V^K$. That is to say, $P$ maps $V$ into $V^K$
and $Pv=v$ for all $v\in V^K$.
\end{lemma}
Here $V^K=\{v\in V:\Pi(x)v=v\text{ for all }x\in K\}$.
\begin{proof}
For $y\in K$, left invariance gives
$\Pi(y)Pv=\int_K\Pi(yx)v\,dx=\int_K\Pi(x)v\,dx=Pv$, so $Pv\in V^K$. If
$v\in V^K$, then $Pv=\int_K\Pi(x)v\,dx=\big(\int_Kdx\big)v=v$ by the
normalization.
\end{proof}
Since $P^2=P$, $P$ is diagonalizable with eigenvalues $0$ and $1$, and
$\operatorname{tr}P=\dim V^K$.
\begin{theorem}[Hall, Theorem~12.15, p.~352 \cite{BCHall2015}]\label{haar:thm:hall-12-15}
If $(\Pi,V)$ and $(\Sigma,W)$ are irreducible representations of $K$, then
\[
 \int_K\overline{\operatorname{trace}(\Pi(x))}\operatorname{trace}(\Sigma(x))\,dx
 =\begin{cases}1&\text{if }V\cong W\\ 0&\text{if }V\not\cong W\end{cases},
\]
where $dx$ is the normalized left-invariant volume form on $K$.
\end{theorem}
\begin{proof}
Hall writes the integrand as the character of $V^*\otimes W$ and identifies
$(V^*\otimes W)^K$ with intertwiners (pp.~353--354). The same argument,
written directly on $\operatorname{Hom}(V,W)$, runs as follows. By the
proof of Theorem~\ref{haar:thm:hall-4-28}, choose inner products for which
$\Pi$ and $\Sigma$ are unitary. Let $K$ act on $\operatorname{Hom}(V,W)$ by
\begin{equation}
 \mathcal T(x)A=\Sigma(x)\,A\,\Pi(x)^{-1}.
\end{equation}
Its invariant vectors are the $A$ with $\Sigma(x)A=A\Pi(x)$, the
intertwiners. In bases, the matrix unit $E_{ab}$ is sent to
$\Sigma(x)E_{ab}\Pi(x)^{-1}$, whose $(a,b)$ entry is
$\Sigma(x)_{aa}\,(\Pi(x)^{-1})_{bb}$ (no sum). Summing over $a$ and $b$, and
using $\Pi(x)^{-1}=\Pi(x)^\dagger$,
\begin{equation}
 \operatorname{tr}\mathcal T(x)=\operatorname{tr}\Sigma(x)\,\operatorname{tr}\Pi(x)^{-1}
 =\operatorname{tr}\Sigma(x)\,\overline{\operatorname{tr}\Pi(x)} .
\end{equation}
Integrate and apply Lemma~\ref{haar:lem:hall-12-16} to $\mathcal T$:
\begin{equation}
 \int_K\overline{\chi_\Pi}\,\chi_\Sigma\,dx
 =\operatorname{tr}\int_K\mathcal T(x)\,dx=\operatorname{tr}P
 =\dim\{\text{intertwiners }V\to W\},
\end{equation}
which is $1$ or $0$ by Schur's lemma.
\end{proof}
Hall's Theorem~12.18 (p.~354) adds completeness: a continuous class function
orthogonal to every irreducible character is zero.

\paragraph{The characters of $SU(2)$ (Hall, Examples~4.10 and 12.23).} Let
$V_m$ be the homogeneous polynomials of degree $m$ in $z=(z_1,z_2)$, with
$[\Pi_m(U)f](z)=f(U^{-1}z)$ (Hall, Eq.~(4.2), p.~82); this is a
representation because $\Pi_m(U_1)\Pi_m(U_2)f(z)=f(U_2^{-1}U_1^{-1}z)
=f((U_1U_2)^{-1}z)$. For $U=\operatorname{diag}(e^{i\theta},e^{-i\theta})$,
$U^{-1}z=(e^{-i\theta}z_1,e^{i\theta}z_2)$, so the basis monomial
$z_1^{m-k}z_2^k$ is multiplied by $e^{-i(m-k)\theta}e^{ik\theta}
=e^{i(2k-m)\theta}$. Sum the geometric series, for $\sin\theta\ne0$:
\begin{align}
 \chi_m(\theta)&=\sum_{k=0}^me^{i(2k-m)\theta}
 =e^{-im\theta}\,\frac{e^{2i(m+1)\theta}-1}{e^{2i\theta}-1}
 =\frac{e^{-im\theta}e^{i(m+1)\theta}}{e^{i\theta}}\cdot
   \frac{e^{i(m+1)\theta}-e^{-i(m+1)\theta}}{e^{i\theta}-e^{-i\theta}}\notag\\
 &=\frac{\sin(m+1)\theta}{\sin\theta}.\label{haar:SU2-character}
\end{align}
For $m=1$, $\chi_1=\frac{\sin2\theta}{\sin\theta}=2\cos\theta
=\operatorname{tr}U$. In spin notation $j=m/2$, $\chi_j(t)=\frac{\sin(2j+1)t}{\sin t}$.
By Eq.~\eqref{haar:weyl-SU2},
\begin{align}
 \langle\chi_m,\chi_n\rangle
 &=\frac2\pi\int_0^\pi\frac{\sin(m+1)t\,\sin(n+1)t}{\sin^2t}\,\sin^2t\,dt
 =\frac1\pi\int_0^\pi\big(\cos(m-n)t-\cos(m+n+2)t\big)\,dt\notag\\
 &=\begin{cases}1&m=n\\0&m\ne n\end{cases}.\label{haar:SU2-character-orthonormal}
\end{align}
For $m=n$ the first cosine is $1$ and the second integrates to $0$ since
$m+n+2\ge2$. The norm $1$ shows that $V_m$ is irreducible. By
Theorem~\ref{haar:thm:hall-4-28}, $V_m\cong\bigoplus_in_iW_i$ with distinct
irreducible $W_i$, so $\chi_m=\sum_in_i\chi_{W_i}$, and
Theorem~\ref{haar:thm:hall-12-15} gives
$\langle\chi_m,\chi_m\rangle=\sum_in_i^2$. This equals $1$ only for a single
summand with $n_i=1$. The $V_m$ are pairwise non-isomorphic since their
characters are orthogonal.
